\chapter{Electrical Safety Risk Assessment}
\label{chap:Electrical Safety Risk Assessment}

%---------------------------- SECTION ----------------------------
\section{Why this assessment matters in construction}

Electricity kills without warning. Unlike many hazards that give sensory warning before the point of injury, a live electrical conductor presents no visible, audible, or olfactory indication of its energy level. Contact with a live conductor at 230\,V AC --- the standard single-phase supply voltage on UK construction sites --- causes immediate ventricular fibrillation at current levels above approximately 50\,mA through the chest, and muscular tetany (the inability to release a grip on the conductor) at lower current levels, meaning that the casualty cannot release themselves from the source of injury. Death from cardiac arrest is the typical outcome of sustained contact with mains-voltage electricity.

Construction sites present an electrical hazard environment that is more complex and more dangerous than almost any fixed workplace. Temporary electrical supplies are installed in stages, extended as the project progresses, and subject to damage from plant movement, weather, and construction activity. The Electricity at Work Regulations 1989 require that electrical systems are constructed, maintained, and used to prevent danger; the 17th and 18th editions of BS 7671 (IET Wiring Regulations) provide the technical standard against which temporary and permanent installations are designed and inspected. Yet the construction site environment routinely contains damaged extension leads, overloaded distribution boards, inadequate earthing, and the persistent risk of contact with buried or overhead electrical services that are identified in Chapter~\ref{chap:Excavation and Groundworks Risk Assessment}.

The principal electrical risks on construction sites fall into three categories: contact with buried or overhead live electrical services (addressed in part by Chapter~\ref{chap:Excavation and Groundworks Risk Assessment}); contact with the temporary electrical supply system (damaged cables, inadequate earthing, overloaded distribution boards, use of inappropriate equipment); and contact with permanent electrical systems in existing buildings during refurbishment or alteration (live circuits whose isolation has not been confirmed). All three require specific assessment and control measures; the common failure mode is the assumption that an electrical system is de-energised when it has not been formally isolated and tested.

The behaviour of small contractors and self-employed tradespeople in relation to electrical safety is a persistent enforcement concern for the HSE. Electrically-connected equipment is used on site without prior isolation; junction boxes are opened without confirming dead; the temporary power system is modified without competent authorisation. Each of these actions places the operative and nearby workers at risk of fatal electrocution from an energy source that provides no warning of its danger.

\barquote{``Assume live until you have confirmed dead, locked out, and tested. There is no such thing as a safe shortcut to electrical isolation.''}

\example{Site Reality}{A general contractor is undertaking a commercial fit-out within an occupied office building. The tenant continues to occupy the floors above the fit-out works. A second-fix electrician is routing cables above a suspended ceiling; he encounters a junction box that he believes is associated with the existing Cat5 data network. He opens the junction box without isolation. The box contains a termination for a 230\,V lighting circuit that serves the floor above; the circuit was live because the tenant's floor above is occupied and the lighting is in use. The electrician contacts the live terminals and receives a hand-to-hand shock across the chest. He collapses; CPR is performed by a colleague and he survives but requires hospitalisation. The investigation identifies the absence of a lockout/tagout procedure, the failure to confirm the nature of the circuit before opening the box, and the absence of a permit-to-work for electrical work in areas adjacent to the live tenant floor.}

%---------------------------- SECTION ----------------------------
\section{Legal and professional framework}

The Electricity at Work Regulations 1989 are the primary legislative instrument for electrical safety in all workplaces including construction sites. Regulation 4 requires that electrical systems are of adequate construction for their purpose and maintained so as to prevent danger. Regulation 13 requires that adequate precautions are taken to prevent electrical equipment that has been made dead from becoming live before work is completed. The Management of Health and Safety at Work Regulations 1999 require suitable and sufficient risk assessment of all electrical work activities. The CDM 2015 pre-construction information duties require the identification of existing live electrical installations in the building.

BS 7671 (Requirements for Electrical Installations, IET Wiring Regulations, 18th edition) provides the technical standard for the design, installation, inspection, and testing of fixed and temporary electrical installations. For construction sites specifically, BS 7671 Part 7 (Particular Installations, Section 704 --- Construction and Demolition Site Installations) provides the construction-site-specific requirements including the mandatory use of reduced low voltage (RLV) systems for portable electrical tools.

\paragraph{Key frameworks relevant to this chapter include: Electricity at Work Regulations 1989; BS 7671:2018 Section 704; MHSWR 1999; CDM 2015; HSG141 Electrical Safety on Construction Sites.}

\begin{itemize}
  \bitem{Electricity at Work Regulations 1989}{Impose a comprehensive framework of duties for electrical safety; require that electrical systems are adequate for their purpose, maintained, and used safely; require that work on or near live conductors is prohibited unless it is unreasonable to work dead; require adequate precautions for work on systems made dead; require competency of persons working on electrical systems.}
  \bitem{BS 7671 Section 704}{Specifies construction-site-specific electrical installation requirements including: the mandatory use of 110\,V centre-tapped-to-earth (CTE) reduced low voltage systems for portable tools and temporary lighting; the maximum distance between overcurrent protection and socket outlets; the requirements for circuit isolation, earthing, and residual current device (RCD) protection; and the requirement for periodic inspection and testing of the temporary electrical installation.}
  \bitem{HSG141 Electrical Safety on Construction Sites}{Provides the HSE's definitive guidance on temporary electrical supply management, tool inspection systems, lockout/tagout procedures, overhead line risk management, and the safe management of electrical work in existing buildings; compliance with HSG141 represents the minimum acceptable standard for construction site electrical safety management.}
  \bitem{CDM 2015}{Requires pre-construction information to identify the location and condition of existing electrical installations, including overhead lines; requires the Construction Phase Plan to address the arrangements for managing electrical risks including temporary supply management and isolation procedures.}
\end{itemize}

%---------------------------- SECTION ----------------------------
\section{Conducting the assessment using the five-step method}

\subsection{Step 1 --- Identify hazards}
\paragraph{Electrical hazards on construction sites include: contact with live overhead electricity lines by plant, scaffold, or operatives; contact with live buried electricity cables during excavation (addressed in Chapter~\ref{chap:Excavation and Groundworks Risk Assessment}); electric shock from the temporary site supply through damaged cables, inadequate earthing, or overloaded circuit protection; electric shock from portable tools with damaged insulation or failed connections; electric shock from work on or adjacent to live circuits in existing buildings during refurbishment; electrical fire from overloaded circuits, arcing at loose connections, or damaged insulation; arc flash from electrical switchgear operating under fault conditions; and electrocution from immersion of electrical equipment or cables in water during flooded site conditions.}

\subsection{Step 2 --- Decide who or what is affected}
\paragraph{Persons at risk from electrical hazards include: electricians and electrical operatives working directly on electrical systems; all site operatives using portable electrical tools and equipment; plant operators in proximity to overhead lines; excavation operatives (see Chapter~\ref{chap:Excavation and Groundworks Risk Assessment}); operatives adjacent to electrical panels and switchgear; visitors to the site who may touch or contact electrical equipment; and members of the public adjacent to sites where overhead lines are present. First-aiders and rescuers may themselves be at risk if they approach a person injured by electricity while the supply remains live.}

\subsection{Step 3 --- Evaluate likelihood and consequence}
\paragraph{Contact with mains-voltage electricity typically results in cardiac arrest and death; the consequence is Very High (5 on a 5×5 consequence scale). On construction sites with uncontrolled temporary supply systems, damaged cables, and no formal tool inspection system, the likelihood of contact is High. Inherent risk for uncontrolled electrical work must be rated Very High (25/25). The mandatory use of 110\,V CTE reduced low voltage systems for portable tools significantly reduces the consequences of contact (55\,V to earth rather than 230\,V), reducing the likelihood of fatal shock; combined with formal isolation procedures, RCD protection, and tool inspection, residual risk can be reduced to Low.}

\subsection{Step 4 --- Select and implement controls}
\paragraph{The primary controls for construction site electrical safety are: use of 110\,V CTE reduced low voltage for all portable tools and temporary lighting rather than 230\,V; provision of 30\,mA RCD protection at all socket outlets; formal lockout/tagout procedure for all work on or near electrical systems; use of goal-post or diverter systems for overhead lines adjacent to plant routes; regular tool inspection system (combined inspection and test with colour-coding); competency requirements for all electrical work (Part P, JIB Grade); permit-to-work for any work on or near live circuits that cannot be isolated.}

\subsection{Step 5 --- Record and review}
\paragraph{Electrical safety records must include: the periodic inspection and testing record for the temporary site supply (minimum monthly inspection required by BS 7671 Section 704); the portable appliance test records for all electrical tools on site; the lockout/tagout records for all electrical isolation activities; the overhead line risk assessment; and the qualifications records for all electrical operatives. These records must be maintained on site and available for inspection.}

\example{Five-Step Application}{A mechanical contractor is installing plant and equipment in a new commercial kitchen. The building is connected to a permanent supply but the fit-out is not yet complete. Step 1 identifies: live 3-phase distribution board adjacent to kitchen equipment installation area; portable tools used in proximity to water supply pipework connections; overhead busbar trunking installed by another contractor above the work area. Step 2 identifies: three mechanical operatives; electrical contractor operative commissioning adjacent controls; site manager conducting daily inspection. Step 3 rates inherent electrocution risk at High (20/25) given live distribution adjacent to wet conditions. Step 4 implements: all mechanical operatives using 110\,V CTE tools only; 30\,mA RCD protection at all supply points; physical barrier installed between live 3-phase board and kitchen area; permit-to-work issued for any task within 600\,mm of the distribution board; busbar trunking not energised until kitchen fit-out is complete and formally handed over. Residual risk rated Low (6/25). Step 5: weekly site electrical inspection including cable condition check.}

%---------------------------- SECTION ----------------------------
\section{Applying the hierarchy of controls}

\paragraph{The hierarchy for electrical safety places the technical voltage reduction measures (110\,V CTE) and RCD protection at the engineering control level, supported by formal isolation procedures and competency requirements as administrative controls.}

\begin{itemize}
  \bitem{Elimination}{Avoid the need for electrical energy in hazardous areas: use battery-powered tools in wet areas and confined spaces; use pneumatic tools where electrical supply is unavailable or where the electrical risk environment is severe; design installations that can be energised only after all wet or hazardous work is complete.}
  \bitem{Substitution}{Use 110\,V CTE reduced low voltage rather than 230\,V for all portable tools and temporary lighting; use safety extra-low voltage (SELV) for lighting in confined spaces, crawl spaces, and highly conductive locations (Section 706 of BS 7671).}
  \bitem{Engineering controls}{Residual current devices (RCDs) at 30\,mA at all socket outlets; 110\,V CTE transformer-based supply systems for all portable tool circuits; mechanical key-interlocked isolation for high-voltage switchgear; goal-post warning systems under overhead lines; physical barriers around live electrical equipment.}
  \bitem{Administrative controls}{Formal lockout/tagout procedure for all electrical isolation; permit-to-work for work on or near live circuits; regular portable appliance testing and colour-code inspection system; competency requirements for all electrical operatives (JIB grade card, Part P registration, 18th edition qualification); pre-work isolation confirmation procedure before any access to electrical enclosures.}
  \bitem{PPE}{Insulated gloves (Class 0 minimum for LV work) for electrical operatives; arc flash protective clothing for work on or near HV switchgear; safety boots with insulating soles; insulated tools with dual-element insulation for all LV electrical work.}
\end{itemize}

%---------------------------- SECTION ----------------------------
\section{Cadence of assessment and review triggers}

\paragraph{The electrical safety assessment must be reviewed as the temporary supply is extended and as permanent installations are energised; the transition from temporary to permanent supply is a particularly high-risk period.}

\begin{itemize}
  \bitem{Monthly inspection of temporary supply}{BS 7671 Section 704 requires periodic inspection and testing of construction site temporary installations; monthly inspection is the industry minimum; the inspection must include cable condition, RCD operation, earthing continuity, and distribution board integrity.}
  \bitem{Before energising any new circuit}{Any new section of temporary or permanent supply must be inspected and tested by a competent person before energisation; no circuit should be assumed dead on the basis of switch position alone without confirmed isolation testing.}
  \bitem{After any electrical incident or near-miss}{Any electrical shock, arc event, or near-miss must trigger an immediate suspension of work, an investigation, and a review of the electrical risk assessment before work resumes.}
  \bitem{When site activities change}{Introduction of activities involving water (concreting, plastering, screeding, external works in wet weather) adjacent to electrical supply equipment must trigger a review of the electrical supply arrangement and additional protection measures.}
  \bitem{At the transition from temporary to permanent supply}{The commissioning and energising of the permanent electrical installation, and the transfer of supply from temporary to permanent, must be managed under a specific electrical commissioning plan with defined isolation and testing procedures.}
\end{itemize}

%---------------------------- SECTION ----------------------------
\section{Common failure modes}

\begin{itemize}
  \bitem{230\,V portable tools used instead of 110\,V CTE}{230\,V portable tools brought onto site from hire fleets or operatives' own equipment, used without conversion to the site's 110\,V supply. This is one of the most common electrical safety enforcement findings on UK construction sites.}
  \bitem{Damaged extension cables not replaced}{Extension cables with damaged insulation, exposed conductors, or unreliable plug connections are repaired with tape rather than replaced. The repair conceals the damage without restoring the insulation integrity.}
  \bitem{No lockout/tagout procedure}{Electrical circuits are de-energised by switching off at the distribution board rather than by formal isolation with locked-off means; circuits can be inadvertently re-energised by another operative switching on at the board.}
  \bitem{RCDs not tested}{RCDs are installed but their operation is not tested at commissioning or at the required periodic intervals. An RCD that has not been tested since installation may not operate when required.}
  \bitem{No portable appliance test system}{Electrical tools and equipment are not subject to a regular combined inspection and test; damaged tools remain in service because no system exists to identify and remove them.}
  \bitem{Work on live circuits without permit-to-work}{Operatives open electrical enclosures, access junction boxes, or conduct cable work without confirming isolation under a formal procedure; circuits that are expected to be dead are not confirmed dead before access.}
\end{itemize}

\example{Failure Pattern}{A plastering subcontractor is engaged on a school refurbishment during the summer holidays. The permanent electrical installation in the classrooms is energised but the lighting circuits are switched off at the local switches. A plasterer drilling into a classroom ceiling to fix plasterboard strikes a cable at 180\,mm depth. The cable is part of a concealed lighting ring main that is de-energised at the local switch but not isolated at the distribution board. A colleague in the next classroom switches the lighting on to check a finishing detail; the cable is energised and the plasterer receives a shock from his metal drill bit. The investigation identifies the absence of any isolation procedure, the absence of any check of the ceiling zone before drilling, and the absence of any communication between trades about energisation status.}

%---------------------------- CHAPTER SUMMARY ----------------------------
\section{Chapter Summary}

\begin{table}[H]
\centering
\begin{tabularx}{\textwidth}{>{\raggedright\arraybackslash}p{3.2cm}X}
\toprule
\textbf{Assessment Element} & \textbf{Operational Requirement} \\
\midrule
Legal/Professional Basis & Electricity at Work Regulations 1989; BS 7671 Section 704; HSG141. 110\,V CTE mandatory for portable tools; RCD protection at all outlets; formal isolation procedure for all electrical work. \\
Method & Apply the full five-step process and document inherent/residual risk. \\
Controls & 110\,V CTE reduced low voltage; 30\,mA RCD protection; formal lockout/tagout; regular PAT and cable inspection; competency requirements for electrical operatives. \\
Cadence & Monthly supply inspection; before energising new circuits; after any electrical incident; at temporary-to-permanent supply transition. \\
Assurance & Use audit, incident learning, and governance reporting to verify effectiveness. \\
\bottomrule
\end{tabularx}
\end{table}

\begin{itemize}
  \bitem{Key takeaway 1}{110\,V centre-tapped-to-earth reduced low voltage is mandatory for all portable tools on UK construction sites; 230\,V portable tools are a common enforcement finding and a preventable cause of electrocution.}
  \bitem{Key takeaway 2}{Formal lockout/tagout is the only reliable means of ensuring that an electrical circuit will not be inadvertently re-energised during work; switching off at a local switch is not a substitute for confirmed isolation.}
  \bitem{Key takeaway 3}{Monthly inspection and testing of the temporary site electrical supply is a minimum regulatory requirement under BS 7671 Section 704; untested RCDs provide a false assurance of protection.}
  \bitem{Key takeaway 4}{Competency for all electrical work must be demonstrated by current JIB grade card, Part P registration, or equivalent qualification; unlicensed electrical work by incompetent persons is one of the most common serious regulatory non-compliances in the construction industry.}
\end{itemize}

\barquote{In Chapter~\ref{chap:Gas Safety and Hot Work Risk Assessment}, we turn from this assessment to the next domain in the Prudentia framework.}
